\documentclass{article}

% Language setting
% Replace `english' with e.g. `spanish' to change the document language
\usepackage[english]{babel}

% Set page size and margins
% Replace `letterpaper' with `a4paper' for UK/EU standard size
\usepackage[letterpaper,top=2cm,bottom=2cm,left=3cm,right=3cm,marginparwidth=1.75cm]{geometry}

% Useful packages
\usepackage{amsmath}
\usepackage{graphicx}
\usepackage{amssymb}
\usepackage[colorlinks=true, allcolors=blue]{hyperref}
\begin{document}
\section{Problem setup and motivation}

We consider kernel ridge regression (KRR) with a translation-invariant kernel
\[
k(x,x') = k(x-x'), \qquad x,x' \in D \subset \mathbb{R}^d,
\]
where \(d\) is small and the dataset size \(N\) is extremely large, potentially on the order of \(10^9\).
Given training data \(\{(x_i,y_i)\}_{i=1}^N\), the KRR objective is
\[
\min_{f \in \mathcal{H}_k}
\frac{1}{N}\sum_{i=1}^N \bigl(f(x_i)-y_i\bigr)^2 + \lambda \|f\|_{\mathcal{H}_k}^2.
\]
By the representer theorem, the solution has the form
\[
f(x)=\sum_{i=1}^N \alpha_i k(x,x_i),
\]
and the coefficient vector \(\alpha\in\mathbb{R}^N\) solves
\[
(K+\lambda I)\alpha = y,
\]
where \(K_{ij}=k(x_i,x_j)\).



For billion-scale data, the direct function-space formulation above is not practical. Even iterative methods become difficult because:
\begin{enumerate}
    \item the kernel matrix \(K\) is dense and cannot be formed explicitly;
    \item each matrix-vector multiplication with \(K\) is expensive unless special structure is exploited;
    \item even if matrix-vector multiplication is accelerated, the iteration count may remain large due to poor conditioning of \(K\).
\end{enumerate}

Our goal is therefore to design a method that simultaneously achieves:
\begin{enumerate}
    \item \emph{\(N\)-independent per-iteration cost}, after one preprocessing pass through the data;
    \item \emph{reduced iteration count}, by flattening the top spectrum of the linear system;
    \item \emph{controlled approximation error}, relative to the original KRR solution.
\end{enumerate}

The proposed strategy combines two ideas:
\begin{itemize}
    \item the \emph{EFGP} Fourier weight-space formulation for translation-invariant kernels, which converts the original \(N\times N\) system into a structured \(M\times M\) system with FFT-accelerated matrix-vector products;
    \item the \emph{EigenPro} spectral preconditioning principle, which rescales the top eigendirections to reduce the effective condition number and enlarge the stable step size so reduce iteration count.
\end{itemize}

\section{Fourier weight-space surrogate}

\subsection{EFGP approximation}

For translation-invariant kernels, one has the inverse Fourier representation
\[
k(z)=\int_{\mathbb{R}^d} \widehat{k}(\xi)e^{2\pi i \langle \xi,z\rangle}\,d\xi.
\]
The EFGP idea is to discretize this integral on an equispaced frequency grid
\[
J_m=\{-m,-m+1,\dots,m\}^d,
\qquad \xi_j = h j,
\]
yielding the approximation
\[
\widetilde{k}(x-x')
=
\sum_{j\in J_m} h^d \widehat{k}(hj)e^{2\pi i h \langle j,x-x'\rangle}.
\]
Equivalently,
\[
\widetilde{k}(x,x')
=
\sum_{j\in J_m} \phi_j(x)\overline{\phi_j(x')},
\qquad
\phi_j(x)=\sqrt{h^d \widehat{k}(hj)}\,e^{2\pi i h \langle j,x\rangle}.
\]
Hence the approximate kernel matrix factorizes as
\[
\widetilde{K}=\Phi\Phi^*, \Phi \in \mathbb{C}^{N\times M}
\]
where
\[
\Phi_{n,j}=\phi_j(x_n),
\qquad
M=(2m+1)^d.
\]
So under this setting, \(d\) should not be too large

\subsection{Weight-space KRR system}

Replacing \(K\) by \(\widetilde{K}=\Phi\Phi^*\), the approximate KRR system becomes
\[
(\widetilde{K}+\lambda I)\widetilde{\alpha}=y.
\]
Using the standard dual-to-weight-space transformation, define
\[
\beta=\Phi^* \widetilde{\alpha}.
\]
Then \(\beta\in\mathbb{C}^M\) solves the weight-space system
\[
A\beta=b, 
\qquad
A:=\Phi^*\Phi+\lambda I,
\qquad
b:=\Phi^*y,
\qquad (\Phi^*\Phi+\lambda I)\beta=\Phi^*y
\]
Once \(\beta\) is computed, the predictor is
\[
\widetilde{f}(x)=\sum_{j\in J_m} \beta_j \phi_j(x)=\phi(x)^*\beta.
\]


\section{Structured matrix-vector multiplication}

\subsection{Toeplitz structure}

Write
\[
\Phi=F D,
\qquad
D_{jj}=\sqrt{h^d\widehat{k}(hj)},
\qquad
F_{n,j}=e^{2\pi i h\langle j,x_n\rangle}.
\]
Then
\[
A = D(F^*F)D + \lambda I.
\]
A key observation is that
\[
(F^*F)_{j,j'}
=
\sum_{n=1}^N e^{2\pi i h\langle j-j',x_n\rangle},
\]
so this entry depends only on the difference \(j-j'\). Thus \(F^*F\) is a \(d\)-dimensional Toeplitz matrix.


Define
\[
v_\ell=\sum_{n=1}^N e^{2\pi i h\langle \ell,x_n\rangle},
\qquad
\ell\in J_{2m}.
\]
Then
\[
[(F^*F)u]_j = \sum_{j'\in J_m} v_{j-j'}u_{j'},
\]
which is a non-periodic discrete convolution.

\subsection{Precomputation via NUFFT}

The convolution kernel \(v\) and the right-hand side \(b=\Phi^*y\) are computed using type-1 NUFFT:
\[
v_\ell=\sum_{n=1}^N e^{2\pi i h\langle \ell,x_n\rangle},
\qquad
b_j=\sum_{n=1}^N \overline{\phi_j(x_n)}\, y_n.
\]
This costs
\[
O(N+M\log M)
\]
and is performed only once.

\section{Why EFGP alone is not enough}

Although EFGP removes the dependence of per-iteration cost on \(N\), it does not automatically solve the conditioning problem. Indeed, those papers show that when the Fourier surrogate is accurate, the weight-space condition number is still comparable to that of the original system with sharp bound
\[
\kappa(A) \leq N/\lambda+1
\]
in the low-noise or small-regularization regime. Thus, EFGP addresses the \emph{cost per iteration}, but not the \emph{number of iterations}. For billion-scale data, this remaining bottleneck becomes dominant.

\section{EigenPro-style spectral preconditioning in weight space}
\subsection{Motivation from EigenPro}

EigenPro was originally developed for kernel interpolation / kernel machines, where mini-batch SGD saturates at a critical batch size
\[
m^*(k)\approx \frac{\beta(K)}{\lambda_1(K)},
\]
with \(\lambda_1(K)\) the largest eigenvalue and \(\beta(K)=\max_i k(x_i,x_i)\). (For mathematical simplicity, we leave aside the problem of parallel capacity and resource memory, which can be checked in[Ma 2019]). The key idea is to flatten the top spectrum by rescaling the first few eigendirections. This increases the stable step size and reduces the number of iterations.


In our setting, maybe the exact same spectral principle can be applied to the EFGP weight-space system. The only difference is that the matrix being preconditioned is now
\[
A=\Phi^*\Phi+\lambda I
\]
rather than the original kernel matrix \(K\).


\subsection{Top-eigenspace approximation}
Suppose we can estimate top eigenpairs
\[
(\theta_i,u_i),\qquad i=1,\dots,q,
\]
for the weight-space matrix \(A\). I need to check the code and methods later.

\subsection{Weight-space EigenPro preconditioner}

Let
\[
U_q=[u_1,\dots,u_q],\qquad
\Theta_q=\mathrm{diag}(\theta_1,\dots,\theta_q),
\qquad
\mu\approx \theta_{q+1}.
\]
Define the spectral preconditioner
\[
P_q
=
I-U_q\Bigl(I-\mu \Theta_q^{-1}\Bigr)U_q^*.
\]
Equivalently,
\[
P_q u
=
u-\sum_{i=1}^q \left(1-\frac{\mu}{\theta_i}\right)\langle u_i,u\rangle u_i.
\]
This is the direct finite-dimensional analogue of the EigenPro operator preconditioner.

\section{Proposed algorithm}

\subsection{Algorithmic idea}

The overall algorithm has two stages:
\begin{enumerate}
    \item \textbf{EFGP stage:} build a structured Fourier weight-space surrogate so that all matrix-vector products cost \(O(M\log M)\) and no longer depend on \(N\);
    \item \textbf{EigenPro stage:} estimate the dominant weight-space eigenspace and flatten the top spectrum to reduce iteration count.
\end{enumerate}

\subsection{Algorithm outline}

\paragraph{Input.}
Training data \(\{(x_i,y_i)\}_{i=1}^N\) with \(x_i\in\mathbb{R}^d\), a translation-invariant kernel \(k\), regularization \(\lambda>0\), Fourier tolerance \(\varepsilon\), top-eigenspace rank \(q\), and stopping tolerance \(\rho\).

\paragraph{Output.}
Approximate predictor
\[
\widetilde{f}(x)=\sum_{j\in J_m}\beta_j \phi_j(x).
\]

\begin{enumerate}
    \item \textbf{Choose Fourier discretization parameters \((h,m)\).}  
    Use the error formulas from [barnett2023uniform] to choose the grid spacing \(h>0\) and the frequency cutoff \(m\in\mathbb{N}\). Define
    \[
    J_m:=\{-m,-m+1,\dots,m\}^d,
    \qquad
    M:=|J_m|=(2m+1)^d.
    \]

    \item \textbf{Construct the Fourier basis and weight-space system.}  
    Define the basis functions
    \[
    \phi_j(x)=\sqrt{h^d\widehat{k}(hj)}\,e^{2\pi i h\langle j,x\rangle},
    \qquad j\in J_m.
    \]
    Let the design matrix
    \[
    \Phi\in\mathbb{C}^{N\times M},
    \qquad
    \Phi_{n,j}=\phi_j(x_n),
    \]
    where rows are indexed by \(n=1,\dots,N\) and columns by \(j\in J_m\).  
    The weight-space unknown is
    \[
    \beta\in\mathbb{C}^{M}.
    \]
    Define the system matrix and right-hand side
    \[
    A:=\Phi^*\Phi+\lambda I_M \in \mathbb{C}^{M\times M},
    \qquad
    b:=\Phi^*y\in\mathbb{C}^{M},
    \]
    where \(y=(y_1,\dots,y_N)^\top\in\mathbb{R}^N\) and \(I_M\in\mathbb{R}^{M\times M}\).

    \item \textbf{Precompute the Toeplitz kernel and right-hand side.}  
    Define
    \[
    J_{2m}:=\{-2m,-2m+1,\dots,2m\}^d,
    \qquad
    M_{2}:=|J_{2m}|=(4m+1)^d.
    \]
    Compute the Toeplitz kernel vector
    \[
    v\in\mathbb{C}^{M_2},
    \qquad
    v_\ell=\sum_{n=1}^N e^{2\pi i h\langle \ell,x_n\rangle},
    \qquad \ell\in J_{2m},
    \]
    and the right-hand side
    \[
    b\in\mathbb{C}^M,
    \qquad
    b_j=\sum_{n=1}^N \overline{\phi_j(x_n)}\,y_n,
    \qquad j\in J_m,
    \]
    using type-1 NUFFT.  
    Precompute
    \[
    \widehat v \in \mathbb{C}^{M_2},
    \]
    the \(d\)-dimensional FFT of the zero-padded Toeplitz kernel vector \(v\).

    \item \textbf{Build a fast matrix-vector oracle.}  
    Define the diagonal matrix
    \[
    D\in\mathbb{R}^{M\times M},
    \qquad
    D_{jj}=\sqrt{h^d\widehat{k}(hj)},
    \]
    and the Fourier feature matrix
    \[
    F\in\mathbb{C}^{N\times M},
    \qquad
    F_{n,j}=e^{2\pi i h\langle j,x_n\rangle},
    \]
    so that \(\Phi=FD\).  
    For any vector
    \[
    u\in\mathbb{C}^M,
    \]
    compute
    \[
    Au = D(F^*F)(Du)+\lambda u,
    \]
    where
    \[
    F^*F\in\mathbb{C}^{M\times M}.
    \]
    This multiplication is carried out using two FFTs and diagonal multiplications, in \(O(M\log M)\).

    \item \textbf{Estimate the top eigenspace of \(A\).}  
    Obtain approximate top eigenpairs
    \[
    (\theta_i,u_i),\qquad i=1,\dots,q,
    \]
    where
    \[
    \theta_i\in\mathbb{R},\qquad u_i\in\mathbb{C}^M,\qquad \|u_i\|_2=1.
    \]
    Stack them into
    \[
    U_q:=[u_1,\dots,u_q]\in\mathbb{C}^{M\times q},
    \qquad
    \Theta_q:=\mathrm{diag}(\theta_1,\dots,\theta_q)\in\mathbb{R}^{q\times q}.
    \]

    \item \textbf{Construct the EigenPro-style preconditioner.}  
    Set
    \[
    P_q
    :=
    I_M-U_q\Bigl(I_q-\mu\Theta_q^{-1}\Bigr)U_q^*
    \in\mathbb{C}^{M\times M},
    \qquad
    \mu=\theta_{q+1}\ \text{or an approximation to it},
    \]
    where \(I_q\in\mathbb{R}^{q\times q}\) is the identity matrix.  
    Then, ideally,
    \[
    \mathrm{spec}(P_qA)=\{\mu,\dots,\mu,\theta_{q+1},\dots,\theta_M\},
    \]
    i.e. the first \(q\) eigenvalues are flattened to \(\mu\).

    \item \textbf{Use preconditioning before solving the system.}  
    Two options are natural:
    \begin{itemize}
        \item \emph{Preconditioned Richardson / gradient iteration:}
        \[
        \beta_{t+1}=\beta_t-\eta P_q(A\beta_t-b),
        \qquad
        \beta_t\in\mathbb{C}^M.
        \]
        \item \emph{Preconditioned conjugate gradient:} use \(P_q\in\mathbb{C}^{M\times M}\) as a spectral preconditioner within PCG applied to
        \[
        A\beta=b,
        \qquad
        A\in\mathbb{C}^{M\times M},\ \beta,b\in\mathbb{C}^M.
        \]
    \end{itemize}
    Stop when the relative residual satisfies conditions:

    \item \textbf{Evaluate the predictor.}  
    For target points \(\{x_\ell^*\}_{\ell=1}^{N_{\mathrm{test}}}\), compute
    \[
    \widetilde{f}(x_\ell^*)=\sum_{j\in J_m}\beta_j\phi_j(x_\ell^*),
    \qquad \ell=1,\dots,N_{\mathrm{test}},
    \]
    where
    \[
    \beta\in\mathbb{C}^{M}.
    \]
    This batched evaluation is performed using type-2 NUFFT.
\end{enumerate}

\section{Problem to be solved}
\begin{enumerate}
    \item How to estimate the top-k eigen pairs
    \item Eigenpro 3.0(Subsampling-based Approximation, Delayed Projections)
    \item (Multi-)GPU, mini-batch 
    \item Parameter estimation and selection; stepsize\(\mu\) etc
    \item Small d
\end{enumerate}

\end{document}
